Сначала прочитайте число $N$.

Чтобы сделать запрос, выведите <<? x>> $(1 \le x \le N)$ и прочитайте ответ:

\begin{itemize}
\item \texttt{<}, если загаданное число меньше $x$;
\item \texttt{>}, если загаданное число больше $x$;
\item \texttt{=}, если загаданное число равно $x$.
\end{itemize}

Можно сделать не более $10$ запросов. Получив ответ \texttt{=}, завершите программу: отдельный итоговый ответ выводить не нужно.

Ваша программа получит вердикт <<Неправильный ответ>> в случае:

\begin{itemize}
\item нарушения формата вывода;
\item запроса числа вне промежутка $[1; N]$;
\item превышения количества доступных запросов;
\item завершения программы без полученного ответа \texttt{=}.
\end{itemize}

\textbf{После каждого выведенного запроса обязательно сбрасывайте буфер вывода}, иначе ваша программа получит вердикт <<Idleness Limit Exceeded>>:

\begin{itemize}
\item C++: \texttt{std::endl} при выводе через \texttt{iostream} (\texttt{cout << ... << endl}), \texttt{fflush(stdout)} после вывода через \texttt{stdio} (\texttt{printf});
\item Java: \texttt{System.out.println(...)}, а при выводе через другой поток (\texttt{PrintWriter}, \texttt{BufferedWriter}) --- его \texttt{flush()} после каждой строки;
\item Python: \texttt{print(..., flush=True)}, а ответы интерактора читайте через \texttt{input()}.
\end{itemize}
